\chapter{PHP2Shell: campaña de explotación de octubre}
\label{chap:incidente1}

\section{Contexto y justificación de selección}
Se analiza la campaña de octubre de 2026 que utilizó PHP2Shell,
dos vulnerabilidades del intérprete PHP, para comprometer
servidores y dispositivos de seguridad. Dream reconstruyó una
operación con 25 servidores comprometidos; Unit 42 confirmó las
herramientas y los resultados comunicados~\cite{dream2026php,unit422026php}.
Se selecciona por mostrar cómo un fallo en un componente compartido
puede convertir un dispositivo perimetral en una vía de acceso
a la infraestructura interna.

\section{Desarrollo del incidente y CVE/CVSS}
La operación combinó ejecución remota mediante PHP2Shell con
una escalada de privilegios del núcleo Linux,
CVE-2026-46300. Según Dream, los atacantes identificaron equipos
PAN-OS expuestos, seleccionaron versiones compatibles con su
explotación y comprometieron sus interfaces de gestión.
Los días 23 y 24 de octubre accedieron a dispositivos de un
entorno que los investigadores consideraron vinculado al
Gobierno estadounidense. Desde un cortafuegos establecieron
un túnel hacia la red interna y comprometieron un servidor
de gestión, obteniendo credenciales y enumerando
Active Directory~\cite{dream2026php}.

PHP2Shell comprende \textbf{CVE-2026-109618}, donde el agotamiento
de memoria deja datos de petición incorrectamente inicializados,
y \textbf{CVE-2026-109619}, donde otra petición atendida por el mismo
trabajador utiliza referencias a memoria liberada
(\textit{use-after-free}). La ejecución inicial tiene los permisos
del trabajador web; el acceso como administrador deriva del
encadenamiento con el fallo del núcleo~\cite{php2026oom,php2026stale,dream2026php}.

PHP Group publica para ambos \textbf{CVSS v4.0: 9,2, crítico},
con el vector base:\par
\noindent{\small\ttfamily
    CVSS:4.0/AV:N/AC:H/AT:P/PR:N/UI:N/\allowbreak
    VC:H/VI:H/VA:H/SC:N/SI:N/SA:N}\par
La valoración refleja acceso por red sin autenticación ni
interacción, complejidad alta y condiciones de ejecución
necesarias. No es un cálculo propio~\cite{php2026oom,php2026stale}.

\section{Impacto y mitigación}
Los compromisos permitieron acceso persistente, extracción de
credenciales y reconocimiento interno; no se acreditó el éxito
de todas las acciones posteriores~\cite{dream2026php}.
Para las organizaciones, esto compromete el control de sus
redes; para quienes utilizan sus servicios, supone riesgo de
exposición de información e interrupción del acceso.
No existe en las fuentes consultadas un balance global de
personas afectadas. Sí hubo consecuencias preventivas observables:
el NCSC neerlandés retiró temporalmente su web y comunicó
cierres similares en otras entidades~\cite{ncsc2026php}.

La corrección exige actualizar a PHP 8.2.35, 8.3.36, 8.4.27
o 8.5.12, según la rama, o aplicar los parches del proveedor.
Las modificaciones de configuración no ofrecen protección
completa~\cite{php2026aviso}. En PAN-OS deben instalarse sus
actualizaciones específicas y restringirse la gestión a redes
de confianza~\cite{unit422026php}.

\section{Análisis crítico}
Amazon observó intentos desde el 17 de octubre, sin confirmar
entonces su éxito~\cite{aws2026php}; CISA incorporó ambos CVE
al catálogo de vulnerabilidades explotadas el
día 26~\cite{cisa2026kev}.
Esta evolución obliga a distinguir exposición, intento y
compromiso. La principal lección es que parchear una
vulnerabilidad no elimina los accesos persistentes ni las
credenciales robadas: la recuperación debe investigar también
los sistemas alcanzados desde el dispositivo inicial.